\newcommand{\titulus}{\nomenFesti{S. Dominici, Presbyteri.}
\dies{Die 8. Augusti.}}
\newcommand{\sineobmv}{Sine officium B.M.V.}
\newcommand{\oratio}{\pars{Oratio.}

\noindent Adiuvet Ecclésiam tuam, Dómine, beátus Domínicus méritis et doctrínis atque pro nobis efficiátur piíssimus intervéntor, qui tuæ veritátis éxstitit prædicátor exímius.

\pars{Pro pace in universo mundo.} \scriptura{Sir. 50, 25; 2 Esdr. 4, 20; \textbf{H416}}

\vspace{-4mm}

\antiphona{II D}{temporalia/ant-dapacemdomine.gtex}

\vfill

\noindent Deus, a quo sancta desidéria, recta consília et iusta sunt ópera: da servis tuis illam, quam mundus dare non potest, pacem; ut et corda nostra mandátis tuis dédita, et hóstium subláta formídine, témpora sint tua protectióne tranquílla.

\noindent Per Dóminum nostrum Iesum Christum, Fílium tuum, qui tecum vivit et regnat in unitáte Spíritus Sancti, Deus, per ómnia sǽcula sæculórum.

\noindent \Rbardot{} Amen.}
\newcommand{\invitatorium}{\pars{Invitatorium.}

\vspace{-4mm}

\antiphona{IV}{temporalia/inv-mirabileminsanctis.gtex}}
\newcommand{\hymnusmatutinum}{\pars{Hymnus}

\cuminitiali{IV}{temporalia/hym-LaetiColentes.gtex}}
\newcommand{\lectioi}{\pars{Lectio I.} \scriptura{Prov. 31, 10-20}

\noindent De libro Proverbiórum.

\noindent Mulíerem fortem quis invéniet? Longe super gemmas prétium eius. Confídit in ea cor viri sui et spóliis non indigébit. Reddet ei bonum et non malum ómnibus diébus vitæ suæ.

\noindent Quæsívit lanam et linum et operáta est delectatióne mánuum suárum. Facta est quasi navis institóris de longe portans panem suum. Et de nocte surréxit dedítque prædam domésticis suis et cibária ancíllis suis. Considerávit agrum et emit eum; de fructu mánuum suárum plantávit víneam. Accínxit fortitúdine lumbos suos et roborávit bráchium suum. Gustávit et vidit quia bona est negotiátio eius; non exstinguétur in nocte lucérna eius. Manum suam misit ad colos, et dígiti eius apprehendérunt fusum. Palmas suas apéruit ínopi et manum suam exténdit ad páuperem.}
\newcommand{\responsoriumi}{\pars{Responsorium 1.} \scriptura{\Rbardot{} Prv. 31, 8; \Vbardot{} ibid. 31, 7; \textbf{H400}}

\vspace{-5mm}

\responsorium{I}{temporalia/resp-verbuminiquum-CROCHU-sinedox.gtex}{}}
\newcommand{\lectioii}{\pars{Lectio II.} \scriptura{Prov. 31, 21-31}

\noindent Non timébit dómui suæ a frigóribus nivis: omnes enim doméstici eius vestíti sunt duplícibus. Stragulátam vestem fecit sibi; byssus et púrpura induméntum eius. Nóbilis in portis vir eius, quando séderit cum senatóribus terræ. Síndonem fecit et véndidit et cíngulum trádidit Chananǽo. Fortitúdo et decor induméntum eius, et ridébit in die novíssimo. Os suum apéruit sapiéntiæ, et lex cleméntiæ in lingua eius. Considerávit sémitas domus suæ et panem otiósa non comédit. Surrexérunt fílii eius et beatíssimam prædicavérunt, vir eius et laudávit eam: «Multæ fíliæ fórtiter operátæ sunt, tu supergréssa es univérsas». Fallax grátia et vana est pulchritúdo; múlier timens Dóminum ipsa laudábitur. Date ei de fructu mánuum suárum, et laudent eam in portis ópera eius.}
\newcommand{\responsoriumii}{\pars{Responsorium 2.} \scriptura{\Rbardot{} Prv. 23, 26 \Vbardot{} ibid. 5, 1; \textbf{H400}}

\vspace{-5mm}

\responsorium{I}{temporalia/resp-praebefili-CROCHU.gtex}{}}
\newcommand{\lectioiii}{\pars{Lectio III.} \scriptura{Libellus de principiis O.P.: Acta canonizationis sancti Dominici: Monumenta O.P. Mist. 16, Romæ 1935, pp. 30 ss., 146-147}

\noindent E váriis scriptis históriæ Ordinis Prædicatórum.

\noindent Tanta morum honestáte Domínicus pollébat, tanto divíni fervóris ímpetu ferebátur, ut ipsum esse vas honóris et grátiæ haud dúbie probarétur. Inerat ei firma valde mentis æquálitas, nisi cum ad compassiónem et misericórdiam turbarétur; et quia cor gaudens exhílarat fáciem, plácidam interióris hóminis compositiónem manifésta de foris benignitáte ac vultus hilaritáte prodébat.

\noindent Ubíque virum evangélicum verbo se exhibébat et ópere. Témpore diúrno cum frátribus sociísve nemo commúnior, nemo iucúndior. Noctúrnis horis nemo vigíliis et obsecratiónibus per omnem modum instántior. Raro loquebátur nisi cum Deo, scílicet orándo, vel de Deo, et de hoc monébat fratres suos.

\noindent Fuit autem ei frequens et speciális quædam ad Deum petítio, ut sibi largíri dignarétur veram caritátem, curándæ et procurándæ salúti hóminum efficácem, árbitrans sese tunc primum fore veráciter membrum Christi, cum se totum pro víribus lucrifaciéndis animábus impénderet, sicut Salvátor ómnium Dóminus Iesus totum se nostram óbtulit in salútem. Et ad hoc opus, alto iam dudum providénte consílio, Fratrum Prædicatórum Ordinem instítuit.

\noindent {\color{gray} Sæpe hortabátur fratres dicti Ordinis, verbis et lítteris suis, quod semper studérent in novo et vétere testaménto. Semper gestábat secum Matthǽi evangélium et epístolas Pauli, et multum studébat in eis, ita quod fere sciébat eas memóriter.}

\noindent {\color{gray} Bis vel ter eléctus fuit in epíscopum et ipse semper rénuit, volens pótius cum frátribus suis in paupertáte vívere quam áliquem episcopátum habére. Virginitátis suæ decus illibátum usque in finem conservávit. Desiderábat flagellári et frustátim incídi et mori pro fide Christi. De quo Gregórius nonus affirmávit: «Novi virum totíus apostólicæ régulæ sectatórem, quem et in cælis non est ambíguum ipsórum Apostolórum glóriæ copulátum».}}
\newcommand{\responsoriumiii}{\pars{Responsorium 3.} \scriptura{\Vbardot{} Sap. 10, 10; \textbf{H379}}

\vspace{-5mm}

\responsorium{VIII}{temporalia/resp-istehomoperfecit-CROCHU-cumdox.gtex}{}}
\newcommand{\hymnuslaudes}{\pars{Hymnus}

\cuminitiali{VIII}{temporalia/hym-NovusAthleta.gtex}}
\newcommand{\lectiobrevis}{\pars{Lectio Brevis.} \scriptura{Heb. 13, 7-9a}

\noindent Mementóte præpositórum vestrórum, qui vobis locúti sunt verbum Dei, quorum intuéntes éxitum conversatiónis imitámini fidem. Iesus Christus heri et hódie idem, et in sǽcula! Doctrínis váriis et peregrínis nolíte abdúci.}
\newcommand{\responsoriumbreve}{\pars{Responsorium breve.} \scriptura{Is. 62, 6}

\cuminitiali{VI}{temporalia/resp-superteierusalem.gtex}}
\newcommand{\benedictus}{\pars{Canticum Zachariæ.}

\vspace{-4mm}

\antiphona{II D}{temporalia/ant-instabatenim.gtex}

\vspace{-2mm}

\scriptura{Lc. 1, 68-79}

\vspace{-2mm}

\cantusSineNeumas
\initiumpsalmi{temporalia/benedictus-initium-ii-D-auto.gtex}

%\vspace{-1.5mm}

\input{temporalia/benedictus-ii-D.tex} \Abardot{}}
\newcommand{\preces}{\noindent Christo, bono pastóri,~\gredagger{} qui pro suis óvibus ánimam pósuit,~\grestar{} laudes grati exsolvámus et supplicémus, dicéntes:

\Rbardot{} Pasce pópulum tuum, Dómine.

\noindent Christe, qui in sanctis pastóribus misericórdiam et dilectiónem tuam dignátus es osténdere,~\grestar{} numquam désinas per eos nobíscum misericórditer ágere.

\Rbardot{} Pasce pópulum tuum, Dómine.

\noindent Qui múnere pastóris animárum fungi per tuos vicários pergis,~\grestar{} ne destíteris nos ipse per rectóres nostros dirígere.

\Rbardot{} Pasce pópulum tuum, Dómine.

\noindent Qui in sanctis tuis, populórum dúcibus, córporum animarúmque médicus exstitísti,~\grestar{} numquam cesses ministérium in nos vitæ et sanctitátis perágere.

\Rbardot{} Pasce pópulum tuum, Dómine.

\noindent Qui, prudéntia et caritáte sanctórum, tuum gregem erudísti,~\grestar{} nos in sanctitáte iúgiter per pastóres nostros ædífica.

\Rbardot{} Pasce pópulum tuum, Dómine.}
\newcommand{\benedicamuslaudes}{\cuminitiali{}{temporalia/benedicamus-memoria-laudes.gtex}}
\include{hebdomadaxviii}
\include{sabbato}
